\documentclass{article}
\usepackage{amsmath, amssymb, amsthm}
\usepackage{hyperref}
\usepackage{geometry}
\geometry{margin=1in}

\title{Erd\H{o}s Problem \#1075}
\date{Last updated: 2025-10-05}
\author{Source: erdosproblems.com}

\begin{document}
\maketitle

\noindent\textbf{Status:} OPEN \quad \textbf{No prize}

\noindent\textbf{Tags:} hypergraphs

\medskip
\noindent\textbf{Problem Statement:}

Let $r\geq 3$. There exists $c_r>r^{-r}$ such that, for any $\epsilon>0$, if $n$ is sufficiently large, the following holds.

Any $r$-uniform hypergraph on $n$ vertices with at least $(1+\epsilon)(n/r)^r$ many edges contains a subgraph on $m$ vertices with at least $c_rm^r$ edges, where $m=m(n)\to \infty$ as $n\to \infty$.

\bigskip
\noindent\textbf{Remarks:}

Erd\H{o}s \cite{Er64f} proved that this is true with $c_r=r^{-r}$ whenever the graph has at least $\epsilon n^r$ many edges.

\bigskip
\noindent\textbf{References:}

[Er64f] Erd\H{o}s, P., On extremal problems of graphs and generalized graphs. Israel J. Math. (1964), 183--190.

\bigskip
\noindent\small{Source: \url{https://www.erdosproblems.com/1075}}
\end{document}
